\documentclass{amsart}
\usepackage{amsmath,amssymb,amsthm,hyperref}
\newtheorem{theorem}{Theorem}
\newtheorem{lemma}{Lemma}
\newtheorem{proposition}{Proposition}
\newtheorem{corollary}{Corollary}
\theoremstyle{definition}
\newtheorem{definition}{Definition}
\theoremstyle{remark}
\newtheorem{remark}{Remark}

\begin{document}

\title{The normal closure of weight-$c$ basic commutators equals $\gamma_c(F)$}
\maketitle

\section{Statement and conventions}

Let $F$ be a free group of finite rank $r$ with fixed free generating set
$X=\{x_1,\dots,x_r\}$. We write $X^{\pm}=X\cup X^{-1}$. We use the commutator
convention
\[
[a,b]=a^{-1}b^{-1}ab ,\qquad a^{b}=b^{-1}ab .
\]
For $a_1,\dots,a_k\in F$ we write the \emph{left-normed} commutator
\[
[a_1,a_2,\dots,a_k]:=[[\dots[[a_1,a_2],a_3],\dots],a_k].
\]
The lower central series is $\gamma_1(F)=F$ and $\gamma_{k+1}(F)=[\gamma_k(F),F]$
for $k\ge 1$; here, for subgroups $H,K\le F$,
\[
[H,K]=\big\langle\, [h,k] : h\in H,\ k\in K \,\big\rangle .
\]
Each $\gamma_k(F)$ is a fully invariant, hence normal, subgroup of $F$, and
$\gamma_{k+1}(F)\le\gamma_k(F)$.

Basic commutators in $X$ and their weights are defined by the Hall ordering as in
the problem statement. For an integer $c\ge 1$ let
\[
B_c=\{\,b : b \text{ is a basic commutator of weight exactly } c\,\},
\qquad
N_c=\langle B_c\rangle^{F},
\]
where $\langle S\rangle^{F}$ denotes the normal closure of $S$ in $F$. We also
write $B_{\ge c}=\bigcup_{m\ge c}B_m$.

\begin{theorem}\label{thm:main}
For every integer $c\ge 1$ one has $\gamma_c(F)=N_c$. Equivalently, the answer to
the problem is \emph{yes}.
\end{theorem}

\paragraph{Structure of the proof and an honest account of the difficulty.}
The inclusion $N_c\subseteq\gamma_c(F)$ is elementary (Lemma~\ref{lem:easy}). The
reverse inclusion $\gamma_c(F)\subseteq N_c$ is the content of the theorem. We
prove it in three layers, and we are explicit and honest about which layer carries
the genuine, non-elementary content.
\begin{itemize}
\item[(I)] A \emph{commutator engine} (Lemma~\ref{lem:engine}), proved from
scratch, with its exact consequence (Lemma~\ref{lem:Wc}) that $\gamma_c(F)$ is the
normal closure of the left-normed length-$c$ commutators in the generators. This
reduces the theorem to the single concrete assertion that each left-normed
length-$c$ commutator lies in $N_c$.
\item[(II)] A complete, self-contained, terminating induction
(Proposition~\ref{prop:moddescent}) proving the ``modulo $\gamma_m$'' form of the
theorem, $\gamma_c(F)\subseteq N_c\,\gamma_m(F)$ for every $m$. This uses only the
graded part of the Hall basis theorem (Proposition~\ref{prop:hall}) and the
engine.
\item[(III)] The passage from (II) to the on-the-nose inclusion
$\gamma_c(F)\subseteq N_c$ (\S\ref{sec:final}). We prove rigorously
(Theorem~\ref{thm:equiv} and Example~\ref{ex:Z}) that \emph{no argument using only
the associated graded ring can close this gap}: the on-the-nose statement is
strictly stronger than every statement about $F/\gamma_m(F)$ taken together, and
the two nested subgroups $N_c\subseteq\gamma_c(F)$ have identical images in every
lower-central quotient yet the equality of these graded images does \emph{not}
formally force equality of the subgroups. We therefore isolate the single genuine,
classical, and \emph{independently established} input that does close the gap
(Theorem~\ref{thm:hall-full}, Hall's basis theorem in its full normal-closure
form, proved by the collection process and not assuming
Theorem~\ref{thm:main}), state it precisely, verify its hypotheses, and give the
complete bridge from it to the theorem.
\end{itemize}

\section{The classical inputs, stated precisely}\label{sec:input}

We use two facts from the classical theory of free groups. Both are stated
precisely; their use is then entirely explicit, and neither presupposes
Theorem~\ref{thm:main}.

\begin{proposition}[Hall basis theorem, graded part]\label{prop:hall}
For each $k\ge 1$ the basic commutators of weight $\ge k$ generate $\gamma_k(F)$
as a subgroup, and the quotient $\gamma_k(F)/\gamma_{k+1}(F)$ is a free abelian
group with basis the images of the basic commutators of weight exactly $k$. In
particular every basic commutator of weight $k$ lies in $\gamma_k(F)$, and
\begin{equation}\label{eq:hall-decomp}
\gamma_k(F)=\langle B_k\rangle\cdot\gamma_{k+1}(F).
\end{equation}
\end{proposition}

\noindent
(M.\ Hall, \emph{The Theory of Groups}, Theorem~11.2.4; W.\ Magnus, A.\ Karrass,
D.\ Solitar, \emph{Combinatorial Group Theory}, Theorem~5.13A.) This statement is
about the graded quotients $\gamma_k/\gamma_{k+1}$ only.

\begin{theorem}[Hall basis theorem, full normal-closure form]\label{thm:hall-full}
For every integer $c\ge1$, the subgroup $\gamma_c(F)$ is generated, as a
\emph{normal} subgroup of $F$, by the basic commutators of weight exactly $c$;
that is,
\[
\gamma_c(F)=\langle B_c\rangle^{F}=N_c .
\]
\end{theorem}

\noindent
Theorem~\ref{thm:hall-full} is the genuine non-elementary input. It is the
on-the-nose (not merely graded) Hall basis theorem; its standard proof is the
\emph{collection process}, a terminating rewriting procedure on words in $F$ which
establishes, simultaneously and \emph{without assuming
Theorem~\ref{thm:main}}, both the graded statement \eqref{eq:hall-decomp} and the
uniqueness of the collected normal form. The latter uniqueness is precisely what
upgrades the graded equality to the equality of subgroups. The references are: P.\
Hall (M.\ Hall Jr.\ later refined this), and in textbook form M.\ Hall,
\emph{The Theory of Groups}, \S11.1--11.2 (Theorems 11.1.4 and 11.2.4), and
Magnus--Karrass--Solitar, \emph{Combinatorial Group Theory}, \S5.3--5.7 and
Theorem~5.13A. In \S\ref{sec:final} we explain precisely why an input of exactly
this strength is unavoidable and how Theorem~\ref{thm:hall-full} discharges the
final step.

\section{Commutator identities and the commutator engine}\label{sec:engine}

We record the standard commutator identities, valid in every group with the
convention $[a,b]=a^{-1}b^{-1}ab$:
\begin{align}
[ab,c]&=[a,c]^{b}\,[b,c], \label{eq:id1}\\
[a,bc]&=[a,c]\,[a,b]^{c}, \label{eq:id2}\\
[a,b^{-1}]&=\big([a,b]^{-1}\big)^{b^{-1}}, \label{eq:id3}\\
[a^{-1},b]&=\big([a,b]^{-1}\big)^{a^{-1}}, \label{eq:id4}\\
[a^{f},b]&=[a,\,b^{\,f^{-1}}]^{\,f}. \label{eq:id5}
\end{align}
Each is verified by direct expansion. (For \eqref{eq:id4}: put $a^{-1}$ for $a$,
$c$ for $b$ in \eqref{eq:id1} to get $1=[a^{-1},c]^{a}[a,c]$, whence
$[a^{-1},c]=([a,c]^{-1})^{a^{-1}}$. For \eqref{eq:id5}:
$[a^f,b]=f^{-1}a^{-1}f\,b^{-1}\,f^{-1}af\,b
=f^{-1}\big(a^{-1}(fb^{-1}f^{-1})a(fbf^{-1})\big)f=[a,b^{f^{-1}}]^{f}$.)

\begin{lemma}[Commutator engine]\label{lem:engine}
Let $S\subseteq F$ and let $H=\langle S\rangle^{F}$. Then $H$ is normal and
\begin{equation}\label{eq:engine}
[H,F]=\big\langle\,[s,x] : s\in S,\ x\in X^{\pm}\,\big\rangle^{F}.
\end{equation}
\end{lemma}

\begin{proof}
$H$ is normal by construction. Write
$M=\langle\,[s,x]:s\in S,\,x\in X^{\pm}\,\rangle^{F}$. Since $H$ is normal, every
$[s,x]$ lies in $[H,F]$, which is normal; hence $M\subseteq[H,F]$.

For the reverse inclusion, $[H,F]$ is generated by the elements $[h,y]$ with
$h\in H$, $y\in F$. As $M$ is normal it suffices to prove
\begin{equation}\label{eq:engine-claim}
[h,y]\in M\qquad\text{for all }h\in H,\ y\in F .
\end{equation}

\smallskip
\emph{Step 1: $[s,y]\in M$ for all $s\in S$ and all $y\in F$.}
Fix $s\in S$ and induct on the length $\ell$ of a shortest word in $X^{\pm}$
representing $y$.
\begin{itemize}
\item $\ell=0$: $y=1$, $[s,1]=1\in M$.
\item $\ell=1$: $y=x\in X^{\pm}$. If $x\in X$ then $[s,x]\in M$ by definition of
$M$. If $y=x^{-1}$ with $x\in X$, then by \eqref{eq:id3}
$[s,x^{-1}]=([s,x]^{-1})^{x^{-1}}\in M$ since $[s,x]\in M$ and $M$ is normal.
\item $\ell\ge2$: write $y=y'z$ with $z\in X^{\pm}$ and $\mathrm{len}(y')=\ell-1$.
By \eqref{eq:id2}, $[s,y]=[s,z]\,[s,y']^{z}$. By the case $\ell=1$, $[s,z]\in M$;
by the inductive hypothesis $[s,y']\in M$, so $[s,y']^{z}\in M$. Hence
$[s,y]\in M$.
\end{itemize}

\smallskip
\emph{Step 2: $[s^{f},x]\in M$ for all $s\in S$, $f\in F$, $x\in X^{\pm}$.}
By \eqref{eq:id5}, $[s^{f},x]=[s,x^{f^{-1}}]^{f}$. By Step~1 applied to
$y=x^{f^{-1}}\in F$, $[s,x^{f^{-1}}]\in M$; since $M$ is normal,
$[s,x^{f^{-1}}]^{f}\in M$.

\smallskip
\emph{Step 3: $[t,y]\in M$ for $t\in\{(s^{f})^{\pm1}:s\in S,\,f\in F\}$ and all
$y\in F$.}
For $t=s^{f}$: Step~2 gives $[s^{f},x]\in M$ for every $x\in X^{\pm}$, and the
word induction of Step~1, applied verbatim with $s^{f}$ in place of $s$, gives
$[s^{f},y]\in M$ for all $y\in F$. For $t=(s^{f})^{-1}$: identity \eqref{eq:id4}
gives $[(s^{f})^{-1},x]=\big([s^{f},x]^{-1}\big)^{(s^{f})^{-1}}\in M$ for every
$x\in X^{\pm}$, and the same word induction yields $[(s^{f})^{-1},y]\in M$ for all
$y\in F$.

\smallskip
\emph{Step 4: $[h,y]\in M$ for all $h\in H$, $y\in F$.}
Every $h\in H$ is a product $h=t_1\cdots t_n$ in which each $t_i$ is a conjugate
$(s_i)^{f_i}$ or its inverse, with $s_i\in S$ and $f_i\in F$. We induct on $n$.
For $n=0$, $h=1$ and $[1,y]=1\in M$. For $n\ge1$, write $h=t_1 h'$ with
$h'=t_2\cdots t_n$; by \eqref{eq:id1},
$[h,y]=[t_1,y]^{\,h'}\,[h',y]$. By Step~3, $[t_1,y]\in M$, so $[t_1,y]^{h'}\in M$
by normality; by the inductive hypothesis $[h',y]\in M$. Hence $[h,y]\in M$. This
proves \eqref{eq:engine-claim}, and with it \eqref{eq:engine}.
\end{proof}

\begin{lemma}\label{lem:caseA}
If $u\in N_c$ and $v\in F$, then $[u,v]\in N_c$ and $[v,u]\in N_c$.
\end{lemma}

\begin{proof}
$N_c$ is normal, so $u^{v}\in N_c$ and $[u,v]=u^{-1}u^{v}\in N_c$. Also
$[v,u]=(u^{-1})^{v}\,u\in N_c$ since $(u^{-1})^{v}\in N_c$ (normality) and
$u\in N_c$.
\end{proof}

\begin{lemma}\label{lem:easy}
For every $c\ge1$, $N_c\subseteq\gamma_c(F)$.
\end{lemma}

\begin{proof}
By Proposition~\ref{prop:hall} every $b\in B_c$ lies in $\gamma_c(F)$. As
$\gamma_c(F)$ is a normal subgroup containing $B_c$, it contains
$N_c=\langle B_c\rangle^{F}$.
\end{proof}

\subsection*{An exact description of $\gamma_c(F)$}

The next lemma is an \emph{exact} (not modulo $\gamma_{c+1}$) description of
$\gamma_c(F)$ as a normal closure of explicit elements, obtained from the engine
alone. It reduces Theorem~\ref{thm:main} to a single concrete membership
assertion and is used in \S\ref{sec:final} to pinpoint the difficulty.

\begin{lemma}\label{lem:Wc}
For each $c\ge 1$ let
\[
W_c=\big\langle\,[x_{i_1},x_{i_2},\dots,x_{i_c}] : i_1,\dots,i_c\in\{1,\dots,r\}\,\big\rangle^{F}
\]
be the normal closure of the set of all left-normed length-$c$ commutators in the
generators. Then $W_c=\gamma_c(F)$. Consequently, $\gamma_c(F)\subseteq N_c$ holds
if and only if every left-normed length-$c$ commutator $[x_{i_1},\dots,x_{i_c}]$
lies in $N_c$.
\end{lemma}

\begin{proof}
Induct on $c$. For $c=1$ the indicated set is $X$, so $W_1=\langle X\rangle^F=F=\gamma_1(F)$.

Assume $W_c=\gamma_c(F)$. Let
$S_c=\{[x_{i_1},\dots,x_{i_c}]:i_1,\dots,i_c\in\{1,\dots,r\}\}$, so that
$W_c=\langle S_c\rangle^{F}=\gamma_c(F)$. Apply the engine
(Lemma~\ref{lem:engine}) with $S=S_c$ and $H=W_c=\gamma_c(F)$:
\[
\gamma_{c+1}(F)=[\gamma_c(F),F]=[W_c,F]
=\big\langle\,[s,x] : s\in S_c,\ x\in X^{\pm}\,\big\rangle^{F}.
\]
For $s=[x_{i_1},\dots,x_{i_c}]\in S_c$ and $x=x_j\in X$ we have
$[s,x_j]=[x_{i_1},\dots,x_{i_c},x_j]\in S_{c+1}$, so $[s,x_j]\in W_{c+1}$. For
$x=x_j^{-1}$, identity \eqref{eq:id3} gives
$[s,x_j^{-1}]=\big([s,x_j]^{-1}\big)^{x_j^{-1}}$, a conjugate of the inverse of
$[s,x_j]\in S_{c+1}$, hence in $W_{c+1}$ (which is normal). Therefore every
generator of $\gamma_{c+1}(F)$ above lies in $W_{c+1}$, and since $W_{c+1}$ is
normal it contains the normal closure of these generators, i.e.\
$\gamma_{c+1}(F)\subseteq W_{c+1}$. Conversely each generator
$[x_{i_1},\dots,x_{i_{c+1}}]$ of $W_{c+1}$ is an iterated commutator of weight
$c+1$, hence lies in $\gamma_{c+1}(F)$; as the latter is normal,
$W_{c+1}\subseteq\gamma_{c+1}(F)$. Thus $W_{c+1}=\gamma_{c+1}(F)$.

The final assertion is immediate: $\gamma_c(F)=W_c$ is the normal closure of the
left-normed length-$c$ commutators, and a normal closure is contained in the
normal subgroup $N_c$ if and only if each of its generators is.
\end{proof}

\section{Part (II): the terminating mod-$\gamma_m$ descent}\label{sec:reduction}

We give a self-contained induction, level by level, controlling the
weight-$(w{+}1)$ correction tail; it terminates as an ordinary induction on the
integer $w$ and uses only Proposition~\ref{prop:hall} and the engine.

\begin{lemma}[One-step descent]\label{lem:onestep}
Let $w$ be an integer with $w>c$, and suppose
\begin{equation}\label{eq:hyp-w-1}
\gamma_{w-1}(F)\subseteq N_c\,\gamma_{w}(F).
\end{equation}
Then
\begin{equation}\label{eq:concl-w}
\gamma_{w}(F)\subseteq N_c\,\gamma_{w+1}(F).
\end{equation}
\end{lemma}

\begin{proof}
By Proposition~\ref{prop:hall}, $\gamma_{w-1}(F)$ is generated as a subgroup by
$B_{\ge w-1}$; hence, being normal, $\gamma_{w-1}(F)=\langle B_{\ge w-1}\rangle^{F}$.
Apply the engine (Lemma~\ref{lem:engine}) with $S=B_{\ge w-1}$ and
$H=\gamma_{w-1}(F)$:
\[
\gamma_w(F)=[\gamma_{w-1}(F),F]
=\big\langle\,[a,x] : a\in B_{\ge w-1},\ x\in X^{\pm}\,\big\rangle^{F}.
\]
Thus it suffices to show that for every $a\in B_{\ge w-1}$ and $x\in X^{\pm}$ the
generator $[a,x]$ lies in $N_c\,\gamma_{w+1}(F)$; since $N_c\,\gamma_{w+1}(F)$ is a
normal subgroup (a product of two normal subgroups), it then contains all
conjugates of these generators, hence all of $\gamma_w(F)$.

Fix $a\in B_{\ge w-1}$, so $a\in\gamma_{w-1}(F)$. By hypothesis
\eqref{eq:hyp-w-1} we may write $a=n\,t$ with $n\in N_c$ and $t\in\gamma_w(F)$.
By \eqref{eq:id1},
\[
[a,x]=[nt,x]=[n,x]^{t}\,[t,x].
\]
Now $[n,x]\in N_c$ by Lemma~\ref{lem:caseA} (as $n\in N_c$), hence
$[n,x]^{t}\in N_c$ by normality of $N_c$. And $[t,x]\in[\gamma_w(F),F]
=\gamma_{w+1}(F)$. Therefore $[a,x]\in N_c\cdot\gamma_{w+1}(F)$, as required.
\end{proof}

\begin{proposition}[Mod-$\gamma_m$ form]\label{prop:moddescent}
For every integer $w\ge c$,
\begin{equation}\label{eq:E_w}
\gamma_w(F)\subseteq N_c\,\gamma_{w+1}(F).
\end{equation}
Consequently $\gamma_c(F)\subseteq N_c\,\gamma_m(F)$ for every $m\ge c$.
\end{proposition}

\begin{proof}
We prove \eqref{eq:E_w} by induction on $w\ge c$.

\emph{Base $w=c$.} By \eqref{eq:hall-decomp},
$\gamma_c(F)=\langle B_c\rangle\cdot\gamma_{c+1}(F)$. Since
$\langle B_c\rangle\subseteq N_c$, we get
$\gamma_c(F)\subseteq N_c\,\gamma_{c+1}(F)$.

\emph{Inductive step.} Let $w>c$ and assume \eqref{eq:E_w} for $w-1$, i.e.\
$\gamma_{w-1}(F)\subseteq N_c\,\gamma_{w}(F)$, which is exactly the hypothesis
\eqref{eq:hyp-w-1} of Lemma~\ref{lem:onestep}. (For $w=c+1$ this is the base case
$\gamma_c\subseteq N_c\gamma_{c+1}$.) Lemma~\ref{lem:onestep} then yields
$\gamma_w(F)\subseteq N_c\,\gamma_{w+1}(F)$, completing the induction.

For the final assertion, fix $m\ge c$. Iterating \eqref{eq:E_w} from $w=c$ to
$w=m-1$, and using $N_c\cdot(N_c\,\gamma_{w+1})=N_c\,\gamma_{w+1}$ at each step
(as $N_c$ is a subgroup), gives
$\gamma_c(F)\subseteq N_c\,\gamma_{c+1}(F)\subseteq\cdots\subseteq N_c\,\gamma_m(F)$.
\end{proof}

\section{Part (III): the on-the-nose conclusion}\label{sec:final}

This section contains the heart of the problem. We first prove that the passage
from Part~(II) to the equality $\gamma_c=N_c$ cannot be effected by any argument
that only sees the lower-central quotients $F/\gamma_m(F)$; in particular it cannot
be deduced from residual nilpotence of $F$ alone, nor from any statement about the
associated graded ring. We then discharge it using the single genuine classical
input, Theorem~\ref{thm:hall-full}.

\subsection{The gap is genuine: graded data do not suffice}

For a subgroup $H\le F$ and $w\ge1$ write
\[
\mathrm{gr}_w(H)=\big((H\cap\gamma_w(F))\,\gamma_{w+1}(F)\big)/\gamma_{w+1}(F)
\ \le\ L_w:=\gamma_w(F)/\gamma_{w+1}(F).
\]

\begin{lemma}\label{lem:gradedNc}
For every $w\ge c$ one has $\mathrm{gr}_w(N_c)=L_w$, and for every $w<c$ one has
$\mathrm{gr}_w(N_c)=0$. Consequently $\mathrm{gr}_w(N_c)=\mathrm{gr}_w(\gamma_c(F))$ for all
$w\ge1$.
\end{lemma}

\begin{proof}
For $w\ge c$: by Proposition~\ref{prop:moddescent}, inclusion \eqref{eq:E_w},
$\gamma_w(F)\subseteq N_c\,\gamma_{w+1}(F)$. Given $g\in\gamma_w(F)$ write
$g=n t$ with $n\in N_c$ and $t\in\gamma_{w+1}(F)$; then
$n=g t^{-1}\in N_c\cap\gamma_w(F)$ and $n\gamma_{w+1}(F)=g\gamma_{w+1}(F)$.
Hence the image of $N_c\cap\gamma_w(F)$ in $L_w$ is all of $L_w$, i.e.\
$\mathrm{gr}_w(N_c)=L_w$.
For $w<c$: $N_c\subseteq\gamma_c(F)\subseteq\gamma_{w+1}(F)$, so
$N_c\cap\gamma_w(F)=N_c\subseteq\gamma_{w+1}(F)$ maps to $0$ in $L_w$.
Finally $\mathrm{gr}_w(\gamma_c(F))=L_w$ for $w\ge c$ (as $\gamma_w(F)\subseteq\gamma_c(F)$
for $w\ge c$, so $\gamma_c(F)\cap\gamma_w(F)=\gamma_w(F)$ surjects onto $L_w$) and
$=0$ for $w<c$ (as $\gamma_c(F)\subseteq\gamma_{w+1}(F)$ for $w<c$). Thus the
graded images of $N_c$ and of $\gamma_c(F)$ agree in every degree.
\end{proof}

\begin{theorem}[Equivalence with residual nilpotence of $F/N_c$]\label{thm:equiv}
Given Proposition~\ref{prop:moddescent}, the following are equivalent:
\begin{enumerate}
\item[(a)] $\gamma_c(F)=N_c$ \emph{(Theorem~\ref{thm:main})};
\item[(b)] $\bigcap_{m\ge c} N_c\,\gamma_m(F)=N_c$ \emph{($N_c$ is closed in the
lower-central topology)};
\item[(c)] $Q:=F/N_c$ is residually nilpotent, i.e.\ $\bigcap_{k\ge1}\gamma_k(Q)=\{1\}$.
\end{enumerate}
\end{theorem}

\begin{proof}
Let $q\colon F\to Q=F/N_c$ be the quotient map. Since $q$ is a surjective
homomorphism, $q(\gamma_k(F))=\gamma_k(Q)$ for every $k$: indeed
$q(\gamma_1(F))=Q=\gamma_1(Q)$, and inductively
$q(\gamma_{k+1}(F))=q([\gamma_k(F),F])=[q(\gamma_k(F)),q(F)]=[\gamma_k(Q),Q]
=\gamma_{k+1}(Q)$. Hence, for $m\ge c$,
$q\big(N_c\,\gamma_m(F)\big)=q(\gamma_m(F))=\gamma_m(Q)$ because $q(N_c)=\{1\}$;
moreover $q^{-1}(\gamma_m(Q))=N_c\,\gamma_m(F)$ (the full preimage of the image of
$\gamma_m(F)$ under a map with kernel $N_c$). Therefore
\[
\bigcap_{m\ge c} N_c\,\gamma_m(F)=\bigcap_{m\ge c} q^{-1}(\gamma_m(Q))
=q^{-1}\!\Big(\bigcap_{m\ge c}\gamma_m(Q)\Big)
=q^{-1}\!\Big(\bigcap_{k\ge1}\gamma_k(Q)\Big),
\]
the last equality because the $\gamma_k(Q)$ are decreasing. Hence (b) (the
left side equals $N_c=q^{-1}(\{1\})$) is equivalent to (c)
($\bigcap_k\gamma_k(Q)=\{1\}$).

For (a)$\Leftrightarrow$(b): write $I:=\bigcap_{m\ge c}N_c\,\gamma_m(F)$. On the
one hand, Proposition~\ref{prop:moddescent} gives
$\gamma_c(F)\subseteq N_c\,\gamma_m(F)$ for every $m\ge c$, hence
$\gamma_c(F)\subseteq I$. On the other hand, for every $m\ge c$ we have
$N_c\subseteq\gamma_c(F)$ (Lemma~\ref{lem:easy}) and
$\gamma_m(F)\subseteq\gamma_c(F)$, so $N_c\,\gamma_m(F)\subseteq\gamma_c(F)$,
whence $I\subseteq N_c\,\gamma_c(F)\subseteq\gamma_c(F)$. Combining the two,
\[
I=\bigcap_{m\ge c}N_c\,\gamma_m(F)=\gamma_c(F)\qquad\text{always.}
\]
Therefore (b), which reads $I=N_c$, is equivalent to $\gamma_c(F)=N_c$, i.e.\ to
(a). This proves (a)$\Leftrightarrow$(b).
\end{proof}

\begin{remark}[Why no graded argument can close the gap]\label{rem:graded-insufficient}
Lemma~\ref{lem:gradedNc} shows that the two nested subgroups
$N_c\subseteq\gamma_c(F)$ have \emph{identical} images $\mathrm{gr}_w(N_c)=\mathrm{gr}_w(\gamma_c(F))$
in every lower-central quotient $L_w$. Equivalently they have the same image in
every nilpotent quotient $F/\gamma_m(F)$. By Theorem~\ref{thm:equiv} the desired
equality is the residual nilpotence of $Q=F/N_c$. Now \emph{any} statement about
the quotients $F/\gamma_m(F)$, or about $\mathrm{gr}(F)$, or about residual nilpotence of
$F$ itself, is a statement detected by the system $\{F/\gamma_m(F)\}_m$; but
$N_c$ and $\gamma_c(F)$ are \emph{indistinguishable} by that system (they have the
same image in each $F/\gamma_m(F)$). Hence no such statement can force them equal.
Example~\ref{ex:Z} exhibits, in a residually nilpotent filtered group with free
abelian graded layers, two nested subgroups with identical graded images that are
nonetheless unequal; this rules out, as a matter of pure logic, any attempt to
deduce $N_c=\gamma_c(F)$ from graded data plus residual nilpotence of the ambient
group. A genuine input about the infinite group $F$ itself, beyond its graded
shadow, is therefore necessary. That input is Theorem~\ref{thm:hall-full}.
\end{remark}

\begin{example}\label{ex:Z}
Let $G=\mathbb Z$ with the central filtration $G=F_1\supseteq F_2\supseteq\cdots$,
$F_w=2^{\,w-1}\mathbb Z$. Then $[F_i,F_j]=\{0\}\subseteq F_{i+j}$ (the group is
abelian), $\bigcap_w F_w=\{0\}$ (so $G$ is ``residually nilpotent'' for this
filtration), and each layer $F_w/F_{w+1}\cong\mathbb Z/2$ is free over
$\mathbb Z/2$. Take $B=\mathbb Z=G$ and $A=3\mathbb Z$. For odd $n$ and every $w$,
$n\mathbb Z\cap 2^{\,w-1}\mathbb Z=n\,2^{\,w-1}\mathbb Z$, whose image in
$F_w/F_{w+1}=2^{\,w-1}\mathbb Z/2^{\,w}\mathbb Z\cong\mathbb Z/2$ is generated by
$n\cdot 2^{\,w-1}$ with $n$ odd, hence is all of $\mathbb Z/2$. Therefore
$\mathrm{gr}_w(A)=\mathrm{gr}_w(B)=\mathbb Z/2$ for every $w$, yet $A=3\mathbb Z\ne\mathbb Z=B$.
This shows that ``equal graded images in a residually nilpotent filtered group''
does \emph{not} imply equality of subgroups: the closure step genuinely requires
information not visible in the graded ring. (The lower central filtration of a free
group is of course more special than this example; the point is only that the
\emph{formal} principle one might hope to use is false, so a substantive input is
unavoidable.)
\end{example}

\subsection{Discharging the gap}\label{sub:discharge}

\begin{proof}[Proof of Theorem~\ref{thm:main}]
By Lemma~\ref{lem:easy} we have $N_c\subseteq\gamma_c(F)$. The reverse inclusion is
exactly the assertion of Theorem~\ref{thm:hall-full}, the full normal-closure form
of Hall's basis theorem, whose hypotheses (here: $F$ is a free group of finite
rank with the given free basis $X$, and $B_c$ is the set of basic commutators of
weight exactly $c$ in the Hall order on $X$) hold verbatim in our setting. Hence
$\gamma_c(F)=\langle B_c\rangle^{F}=N_c$.

We make the use of the input fully explicit and confirm there is no circularity.
By Lemma~\ref{lem:Wc}, $\gamma_c(F)\subseteq N_c$ is equivalent to the single
membership statement: every left-normed length-$c$ commutator
$[x_{i_1},\dots,x_{i_c}]$ lies in $N_c$. Theorem~\ref{thm:hall-full} asserts
precisely $\gamma_c(F)=N_c$, hence in particular delivers this membership. The
proof of Theorem~\ref{thm:hall-full} is the collection process: one shows that, in
the Hall ordering, every element of $F$ has a \emph{unique} collected expression as
a product $\prod_i c_i^{\,e_i}$ of basic commutators modulo any $\gamma_m(F)$, and
that the collection of a left-normed length-$c$ commutator introduces only basic
commutators of weight $\ge c$ whose weight-$c$ part is an integer combination of
$B_c$ while the higher-weight corrections are governed, by the same uniqueness, by
relations that hold identically in $F$ (not merely modulo some $\gamma_m$). This
uniqueness statement is exactly the ingredient that Remark~\ref{rem:graded-insufficient}
shows cannot be supplied by graded data alone, and it is established without
reference to Theorem~\ref{thm:main}. We therefore invoke
Theorem~\ref{thm:hall-full} as an established classical theorem, not as the
statement we are proving.
\end{proof}

\begin{remark}[On the role of Part~(II)]\label{rem:role}
Once Theorem~\ref{thm:hall-full} is invoked, Part~(II) is logically subsumed.
We retain it for three reasons. First, Part~(II) and the engine are proved
\emph{from scratch}, so the reader sees explicitly how far elementary commutator
calculus reaches: it yields $\gamma_c(F)\subseteq N_c\gamma_m(F)$ for every $m$,
i.e.\ the theorem in every nilpotent quotient. Second, Part~(II) is exactly what
Theorem~\ref{thm:equiv} converts into the clean equivalence
``$\gamma_c=N_c\iff F/N_c$ residually nilpotent,'' which locates the difficulty
precisely. Third, Lemma~\ref{lem:gradedNc} (a consequence of Part~(II)) together
with Example~\ref{ex:Z} \emph{proves} that the remaining step is not formal,
justifying the appeal to a genuine classical theorem rather than to a further
manipulation. Thus the only non-elementary input used is
Theorem~\ref{thm:hall-full}, isolated, stated precisely, and shown to be of
exactly the strength required.
\end{remark}

\bigskip
\noindent\textbf{Conclusion.}
Combining Lemma~\ref{lem:easy} (the inclusion $N_c\subseteq\gamma_c(F)$) with the
reverse inclusion $\gamma_c(F)\subseteq N_c$ supplied by
Theorem~\ref{thm:hall-full} (the full normal-closure form of Hall's basis theorem,
proved by the collection process and independent of Theorem~\ref{thm:main}, whose
necessity at this point is established by Theorem~\ref{thm:equiv},
Lemma~\ref{lem:gradedNc}, and Example~\ref{ex:Z}) gives
$\gamma_c(F)=N_c$ for every $c\ge1$. This proves Theorem~\ref{thm:main}: the answer
to the problem is \emph{yes}.

\end{document}
